\section{Supplementary results}
\label{app:supplementary}

The tables in this appendix are generated directly from canonical JSON/CSV
evidence by \texttt{build\_w1\_latex\_evidence.py}. Their source and output
hashes are recorded in \texttt{W1\_EVIDENCE\_TABLES\_COMPLETE.json}; displayed
precision does not replace the machine-readable values.

\subsection{All validation seeds and frozen-test checkpoints}

Table~\ref{tab:app-validation} exposes seed variability, capacity and selected
epoch for all 15 candidates. Ten runs select the last permitted epoch, so the
common budget may censor convergence even though it is equal across variants.

\begin{table}[p]
\centering\scriptsize
\caption{All validation runs used for frozen model selection. NLL, ADE and FDE are uncalibrated rollout-macro metrics: each is averaged within rollout and then equally across rollouts.}
\label{tab:app-validation}
\begin{tabular}{@{}llrrrrrr@{}}
\toprule
Variant & Seed & Epoch & Trainable & NLL & ADE & FDE & ms/sample \\
\midrule
B1 & 11 & 20 & 1,034,208 & 1.861 & 0.110 & 0.135 & 16.72 \\
B1 & 23 & 20 & 1,034,208 & 1.860 & 0.114 & 0.146 & 16.93 \\
B1 & 37 & 20 & 1,034,208 & 1.861 & 0.107 & 0.131 & 16.33 \\
B2-M & 11 & 18 & 77,600 & 1.997 & 0.914 & 1.626 & 18.32 \\
B2-M & 23 & 19 & 77,600 & 2.029 & 1.088 & 1.722 & 16.49 \\
B2-M & 37 & 20 & 77,600 & 2.026 & 1.053 & 1.807 & 18.05 \\
B2-D & 11 & 20 & 176,096 & 1.873 & 0.248 & 0.449 & 18.37 \\
B2-D & 23 & 19 & 176,096 & 1.873 & 0.251 & 0.444 & 17.64 \\
B2-D & 37 & 20 & 176,096 & 1.872 & 0.244 & 0.440 & 20.15 \\
T1 & 11 & 19 & 86,688 & 2.003 & 0.945 & 1.592 & 18.69 \\
T1 & 23 & 20 & 86,688 & 2.009 & 0.953 & 1.830 & 16.81 \\
T1 & 37 & 20 & 86,688 & 2.024 & 1.046 & 1.722 & 16.61 \\
T2 & 11 & 20 & 165,728 & 1.879 & 0.276 & 0.472 & 18.54 \\
T2 & 23 & 18 & 165,728 & 1.878 & 0.271 & 0.470 & 17.99 \\
T2 & 37 & 20 & 165,728 & 1.877 & 0.268 & 0.455 & 17.14 \\
\bottomrule
\end{tabular}
\end{table}


Table~\ref{tab:app-test} reports the one-shot test after model and representative
seed were frozen on validation. B0 was evaluated as a post-selection reporting
bridge and did not alter the selected B1 stack.

\begin{table}[h]
\centering\small
\caption{One-shot frozen-test results. Calibration was fitted on validation and was not used for architecture ranking.}
\label{tab:app-test}
\begin{tabular}{@{}llrrrrrr@{}}
\toprule
Variant & Seed & Val. rank & NLL & ADE & FDE & Cal. NLL & ms/sample \\
\midrule
B0 & -- & control & 2.171 & 1.299 & 2.685 & 0.582 & 17.62 \\
B1 & 37 & 1 & 1.857 & 0.106 & 0.129 & -2.069 & 18.51 \\
B2-D & 11 & 2 & 1.873 & 0.226 & 0.383 & -1.455 & 18.33 \\
T2 & 23 & 3 & 1.878 & 0.246 & 0.403 & -1.072 & 19.05 \\
T1 & 23 & 4 & 2.004 & 0.928 & 1.704 & -1.287 & 19.78 \\
B2-M & 37 & 5 & 2.024 & 1.047 & 1.730 & -1.351 & 18.74 \\
\bottomrule
\end{tabular}
\end{table}


\subsection{Complete corrected closed-loop contrasts}

Tables~\ref{tab:app-h3-contrasts} and~\ref{tab:app-h4-contrasts} contain every
primary metric contrast, including descriptive cluster intervals, exact
sign-flip p-values and Holm-adjusted values. No adjusted value is confirmatory;
this is consistent with five paired clusters and is not evidence of
equivalence.

\begin{table}[p]
\centering\scriptsize
\caption{All corrected R3 H3 metric contrasts. Effects are B1 minus B0; intervals are deterministic cluster-bootstrap descriptions and tests use five paired initialisation groups.}
\label{tab:app-h3-contrasts}
\begin{tabular}{@{}L{0.25\linewidth}llrL{0.15\linewidth}rrr@{}}
\toprule
Contrast & Metric & Style & Effect & 95\% CI & $p$ & $p_{\mathrm{Holm}}$ & $n$ \\
\midrule
B1--B0, fixed-aggressive & Completion (s) & assertive & -0.790 & [-1.410, -0.180] & 0.1250 & 1.0000 & 5 \\
B1--B0, fixed-aggressive & Separation (m) & assertive & +0.005 & [-0.011, +0.025] & 0.6875 & 1.0000 & 5 \\
B1--B0, fixed-aggressive & Completion (s) & reactive & -0.010 & [-0.090, +0.120] & 1.0000 & 1.0000 & 5 \\
B1--B0, fixed-aggressive & Separation (m) & reactive & -0.011 & [-0.031, +0.000] & 0.3750 & 1.0000 & 5 \\
B1--B0, fixed-medium & Completion (s) & assertive & +0.230 & [-0.630, +1.170] & 0.6250 & 1.0000 & 5 \\
B1--B0, fixed-medium & Separation (m) & assertive & +0.007 & [-0.000, +0.015] & 0.4375 & 1.0000 & 5 \\
B1--B0, fixed-medium & Completion (s) & reactive & -0.220 & [-0.430, -0.020] & 0.1875 & 1.0000 & 5 \\
B1--B0, fixed-medium & Separation (m) & reactive & -0.008 & [-0.023, +0.001] & 0.3750 & 1.0000 & 5 \\
B1--B0, fixed-conservative & Completion (s) & assertive & -0.520 & [-1.320, +0.090] & 0.4375 & 1.0000 & 5 \\
B1--B0, fixed-conservative & Separation (m) & assertive & -0.026 & [-0.075, +0.000] & 0.5000 & 1.0000 & 5 \\
B1--B0, fixed-conservative & Completion (s) & reactive & +0.490 & [+0.020, +0.960] & 0.2500 & 1.0000 & 5 \\
B1--B0, fixed-conservative & Separation (m) & reactive & -0.006 & [-0.016, +0.005] & 0.3750 & 1.0000 & 5 \\
B1--B0, adaptive & Completion (s) & assertive & +0.240 & [-0.090, +0.600] & 0.3750 & 1.0000 & 5 \\
B1--B0, adaptive & Separation (m) & assertive & -0.007 & [-0.018, +0.000] & 0.3125 & 1.0000 & 5 \\
B1--B0, adaptive & Completion (s) & reactive & -0.180 & [-0.660, +0.360] & 0.6250 & 1.0000 & 5 \\
B1--B0, adaptive & Separation (m) & reactive & +0.001 & [-0.016, +0.021] & 0.8125 & 1.0000 & 5 \\
\bottomrule
\end{tabular}
\end{table}


\begin{table}[p]
\centering\scriptsize
\caption{All corrected R3 H4 metric contrasts. Effects are adaptive minus fixed; Holm adjustment is within each registered three-comparator predictor--style family.}
\label{tab:app-h4-contrasts}
\begin{tabular}{@{}L{0.25\linewidth}llrL{0.15\linewidth}rrr@{}}
\toprule
Contrast & Metric & Style & Effect & 95\% CI & $p$ & $p_{\mathrm{Holm}}$ & $n$ \\
\midrule
B1: adaptive vs fixed-aggressive & Completion (s) & assertive & +0.530 & [+0.010, +1.200] & 0.2500 & 0.7500 & 5 \\
B1: adaptive vs fixed-aggressive & Separation (m) & assertive & -0.015 & [-0.037, +0.011] & 0.3750 & 1.0000 & 5 \\
B1: adaptive vs fixed-medium & Completion (s) & assertive & -0.440 & [-1.170, +0.040] & 0.3750 & 0.7500 & 5 \\
B1: adaptive vs fixed-medium & Separation (m) & assertive & -0.000 & [-0.014, +0.011] & 0.8125 & 1.0000 & 5 \\
B1: adaptive vs fixed-conservative & Completion (s) & assertive & +0.440 & [-0.050, +0.970] & 0.4375 & 0.7500 & 5 \\
B1: adaptive vs fixed-conservative & Separation (m) & assertive & -0.003 & [-0.017, +0.009] & 0.6875 & 1.0000 & 5 \\
B1: adaptive vs fixed-aggressive & Completion (s) & reactive & -0.140 & [-0.490, +0.100] & 0.6875 & 1.0000 & 5 \\
B1: adaptive vs fixed-aggressive & Separation (m) & reactive & +0.001 & [-0.000, +0.001] & 0.5000 & 1.0000 & 5 \\
B1: adaptive vs fixed-medium & Completion (s) & reactive & -0.140 & [-0.520, +0.130] & 0.6250 & 1.0000 & 5 \\
B1: adaptive vs fixed-medium & Separation (m) & reactive & +0.001 & [-0.001, +0.002] & 0.5000 & 1.0000 & 5 \\
B1: adaptive vs fixed-conservative & Completion (s) & reactive & -0.420 & [-0.860, +0.070] & 0.1875 & 0.5625 & 5 \\
B1: adaptive vs fixed-conservative & Separation (m) & reactive & -0.002 & [-0.013, +0.010] & 0.9375 & 1.0000 & 5 \\
B0: adaptive vs fixed-aggressive & Completion (s) & assertive & -0.500 & [-1.220, +0.030] & 0.3125 & 0.9375 & 5 \\
B0: adaptive vs fixed-aggressive & Separation (m) & assertive & -0.003 & [-0.023, +0.016] & 0.8750 & 1.0000 & 5 \\
B0: adaptive vs fixed-medium & Completion (s) & assertive & -0.450 & [-1.170, +0.120] & 0.5000 & 1.0000 & 5 \\
B0: adaptive vs fixed-medium & Separation (m) & assertive & +0.014 & [+0.001, +0.026] & 0.1875 & 0.5625 & 5 \\
B0: adaptive vs fixed-conservative & Completion (s) & assertive & -0.320 & [-0.980, +0.170] & 0.5625 & 1.0000 & 5 \\
B0: adaptive vs fixed-conservative & Separation (m) & assertive & -0.023 & [-0.077, +0.008] & 0.9375 & 1.0000 & 5 \\
B0: adaptive vs fixed-aggressive & Completion (s) & reactive & +0.030 & [-0.250, +0.410] & 0.9375 & 1.0000 & 5 \\
B0: adaptive vs fixed-aggressive & Separation (m) & reactive & -0.011 & [-0.051, +0.016] & 0.8125 & 1.0000 & 5 \\
B0: adaptive vs fixed-medium & Completion (s) & reactive & -0.180 & [-0.520, +0.160] & 0.5000 & 1.0000 & 5 \\
B0: adaptive vs fixed-medium & Separation (m) & reactive & -0.008 & [-0.030, +0.013] & 0.5625 & 1.0000 & 5 \\
B0: adaptive vs fixed-conservative & Completion (s) & reactive & +0.250 & [-0.210, +0.710] & 0.3750 & 1.0000 & 5 \\
B0: adaptive vs fixed-conservative & Separation (m) & reactive & -0.008 & [-0.031, +0.015] & 0.4375 & 1.0000 & 5 \\
\bottomrule
\end{tabular}
\end{table}


Table~\ref{tab:app-h4-dominance} applies the registered joint dominance rule to
the same contrasts. It is the source of the 3/12 H4 count and does not search
over comparators or geometry margins.

\begin{table}[h]
\centering\small
\caption{Registered H4 dominance decisions. Effects are adaptive minus fixed; lower completion and higher separation are preferred.}
\label{tab:app-h4-dominance}
\begin{tabular}{@{}llllrrl@{}}
\toprule
Predictor & Style & Fixed comparator & $\Delta t$ (s) & $\Delta d$ (m) & Binary guards & Decision \\
\midrule
B0 & assertive & fixed\_aggressive & -0.500 & -0.0032 & pass & Does not dominate \\
B0 & assertive & fixed\_conservative & -0.320 & -0.0228 & pass & Does not dominate \\
B0 & assertive & fixed\_medium & -0.450 & +0.0136 & pass & Dominates \\
B0 & reactive & fixed\_aggressive & +0.030 & -0.0107 & pass & Does not dominate \\
B0 & reactive & fixed\_conservative & +0.250 & -0.0081 & pass & Does not dominate \\
B0 & reactive & fixed\_medium & -0.180 & -0.0084 & pass & Does not dominate \\
B1 & assertive & fixed\_aggressive & +0.530 & -0.0151 & pass & Does not dominate \\
B1 & assertive & fixed\_conservative & +0.440 & -0.0033 & pass & Does not dominate \\
B1 & assertive & fixed\_medium & -0.440 & -0.0003 & pass & Does not dominate \\
B1 & reactive & fixed\_aggressive & -0.140 & +0.0005 & pass & Dominates \\
B1 & reactive & fixed\_conservative & -0.420 & -0.0019 & pass & Does not dominate \\
B1 & reactive & fixed\_medium & -0.140 & +0.0006 & pass & Dominates \\
\bottomrule
\end{tabular}
\end{table}


\subsection{Calibration and context ablations}

The following generated tables separate three distinct mechanism questions:
whether the Transformer uses its sequence, which B1 input drives aggregate
accuracy, and whether aggregate calibration represents the response-active
tail. An ablation delta establishes sensitivity to an input transformation,
not semantic understanding or closed-loop causality.

\begin{table}[h]
\centering\small
\caption{Post-selection sequence ablations on the frozen test set. Positive deltas mean ablation worsened the metric.}
\label{tab:app-sequence-ablation}
\begin{tabular}{@{}llrrr@{}}
\toprule
Variant & Ablation & $\Delta$NLL & $\Delta$ADE (m) & $\Delta$FDE (m) \\
\midrule
T1 & zero & +0.022 & -0.033 & +0.021 \\
T1 & shuffle & +0.085 & +0.144 & +0.284 \\
T2 & zero & +0.091 & +0.198 & +0.292 \\
T2 & shuffle & +0.149 & +0.350 & +0.428 \\
\bottomrule
\end{tabular}
\end{table}

\begin{table}[h]
\centering\small
\caption{B1 input-shuffle diagnostics across three deterministic mappings.}
\label{tab:app-b1-inputs}
\begin{tabular}{@{}llrrr@{}}
\toprule
Input & Mapping seed & $\Delta$NLL & $\Delta$ADE (m) & Active-tail $\Delta$ADE (m) \\
\midrule
Raster & 20260808 & +0.095 & +0.28883 & +1.907 \\
Target history & 20260808 & -0.000 & +0.00006 & +0.026 \\
Raster & 20260809 & +0.097 & +0.29293 & +1.923 \\
Target history & 20260809 & +0.000 & +0.00019 & +0.029 \\
Raster & 20260810 & +0.086 & +0.27053 & +1.890 \\
Target history & 20260810 & -0.000 & -0.00002 & +0.022 \\
\bottomrule
\end{tabular}
\end{table}

\begin{table}[h]
\centering\small
\caption{Aggregate and response-active calibration diagnostic for the frozen B0/B1 test stacks. The active tail has 15 windows from six rollouts and three initialisation groups.}
\label{tab:app-calibration-tail}
\begin{tabular}{@{}lllrrrr@{}}
\toprule
Subset & Stack & Windows & Rollouts & Init groups & Cal. NLL & Coverage MAE \\
\midrule
all & B0 & 315 & 20 & 5 & 0.582 & 0.407 \\
all & B1 & 315 & 20 & 5 & -2.069 & 0.077 \\
response\_active & B0 & 15 & 6 & 3 & 2.959 & 0.418 \\
response\_active & B1 & 15 & 6 & 3 & 8.573 & 0.443 \\
\bottomrule
\end{tabular}
\end{table}


\subsection{Collision-filtered sensitivity}

Six reactive training rollouts contained 253 raw callbacks. Canonical actor-pair
deduplication and contiguous-frame reconstruction produced 16 contact episodes
over 62 unique frames, primarily target--infrastructure contact. Exact sample-
clock alignment could not be recovered, so the sensitivity conservatively
removed all 162 usable windows from those six rollouts, regenerated train-only
normalisation and reran all 15 validation configurations. The five-model order
remained B1, B2-D, T2, T1, B2-M. This analysis never accessed test labels,
did not change the frozen primary dataset or selected checkpoint, and is
therefore a robustness check rather than a second selection experiment.

\subsection{Evidence not pooled with R3}

Day10--13 analyses use earlier control or geometry implementations and support
diagnostic mechanism statements only. They are absent from
Tables~\ref{tab:app-h3-contrasts}--\ref{tab:app-h4-dominance}, never enlarge
the five-cluster sample and never enter H3/H4 estimates.

\clearpage
